\section{问题二：缺失模态鲁棒情感预测}

\subsection{任务定义}
附件 2 按训练集、验证集和测试集组织。训练集用于学习参数，验证集用于选择最佳 epoch、损失配置和分类策略，测试集只在方案锁定后进行一次最终评估。标准化统计量只由训练集计算。附件 3 没有标签，因此只进行专项推理。

\subsection{DRM-Net 模型}
对每个模态 $m$ 的输入序列 $X_m$ 先进行线性投影，再由 Bi-GRU 得到时序表示 $H_m$。缺失位置由 $M_m$ 标记。跨模态注意力只对有效 key/value 计算归一化权重；当某模态全部缺失时，注意力输出为零向量，避免全缺失输入产生伪均匀注意力。

对有效时间步进行 mask-aware pooling：
\begin{equation}
z_m=\frac{\sum_t M_m(t)H_m(t)}{\max(1,\sum_tM_m(t))}.
\end{equation}
动态门控产生 $g_t,g_a,g_v$，并满足 $g_t+g_a+g_v=1$。融合表示为
\begin{equation}
z=g_tz_t+g_az_a+g_vz_v.
\end{equation}
分类头输出三类极性概率，回归头输出连续情感强度。

\subsection{缺失建模与训练}
本文区分三种情况：原始数据中的真实缺失、训练阶段随机模拟的局部连续缺失，以及附件 3 中的真实专项缺失。训练增强在统一时间轴上随机选择连续区间，将对应模态的输入和 mask 同步置为缺失。缺失类型包括文本、音频、视觉及双模态组合，缺失位置包括前段、中段和后段。

训练损失采用加权交叉熵与 Smooth L1 回归损失的加权和。结构消融分别去除跨模态注意力、动态门控、缺失增强和缺失 mask，保持数据划分、训练轮数和三个 seed 不变。

\subsection{最终评估策略}
三个 seed 分别保存验证集最佳 checkpoint。最终分类概率先求平均，再采用 Argmax：
\begin{equation}
\widehat c=\arg\max_c\frac{1}{3}\sum_{s\in\{42,123,777\}}p_s(c).
\end{equation}
最终测试指标为：Accuracy $66.30\%$，Macro-F1 $62.54\%$，MAE $0.6346$，Pearson $r=0.6630$。普通交叉熵和 Focal Loss 仅作为验证阶段比较方案，不根据测试集选择损失。

\begin{table}[htbp]
\centering
\caption{最终模型与损失函数对照}
\begin{tabular}{lcccc}
\hline
方案 & Accuracy & Macro-F1 & MAE & Pearson $r$ \\ \hline
加权交叉熵（最终） & 66.30\% & 62.54\% & 0.6346 & 0.6630 \\ 
普通交叉熵 & 67.54\% & 59.33\% & 0.6344 & 0.6675 \\ 
Focal Loss & 67.95\% & 60.96\% & 0.6284 & 0.6709 \\ \hline
\end{tabular}
\end{table}

\subsection{结构消融}
\begin{table}[htbp]
\centering
\caption{结构消融结果}
\begin{tabular}{lcccc}
\hline
模型 & Accuracy & Macro-F1 & MAE & Pearson $r$ \\ \hline
Full DRM-Net & 66.30\% & 62.54\% & 0.6346 & 0.6630 \\ 
No Cross-Attention & 66.16\% & 62.17\% & 0.6231 & 0.6723 \\ 
No Dynamic Gate & 66.16\% & 62.69\% & 0.6343 & 0.6597 \\ 
No Missing Augmentation & 64.65\% & 60.83\% & 0.6360 & 0.6646 \\ 
No Missing Mask & 65.75\% & 60.93\% & 0.6428 & 0.6541 \\ \hline
\end{tabular}
\end{table}

缺失增强和缺失 mask 的去除带来更明显的 Macro-F1 下降，支持二者对缺失场景建模的作用。该结论仅针对已运行的实验设置，不外推到未测试的缺失比例或数据分布。

\subsection{附件 3 专项推理}
附件 3 共输出两套各 30 条记录，字段包括样本编号、缺失模态、缺失区间、预测极性、预测强度、不确定度和三模态门控权重。由于附件 3 无标签，本文不计算准确率、F1、MAE 或 Pearson。

